\documentclass[10pt,a4paper]{amsart}
\usepackage[margin=19mm]{geometry}
\usepackage{amsmath,amssymb,mathtools}
\usepackage[scr=rsfs,cal=euler]{mathalfa}
\usepackage{xfrac}
\usepackage[unicode,hidelinks,bookmarks=false]{hyperref}
\newcommand{\ie}{i.e.\ }
\newcommand{\cf}{cf.\ }
\newcommand{\resp}{respectively\ }
\newcommand{\Subsection}{\subsection}
\providecommand{\nobreakdash}{\nobreak\hbox{-}}
\setlength{\emergencystretch}{2em}
\multlinegap0pt
\usepackage{bm}
\usepackage{bbm}
\usepackage{dsfont}
\usepackage{xspace}
\usepackage{cancel}
\usepackage{mathtools}
\usepackage{amsthm}
\makeatletter
\@ifundefined{theorem}{\newtheorem{theorem}{Theorem}}{}
\@ifundefined{coro}{\newtheorem{coro}[theorem]{Corollary}}{}
\makeatother
\DeclareMathOperator{\Tr}{Tr}
\newcommand{\C}{\mathbb{C}}
\newcommand{\II}{\mathbb{I}}
\title[Constructive discretization and approximation in reproducing]{Constructive discretization and approximation in reproducing kernel Hilbert spaces}
\author{Abdellah Chkifa, Matthieu Dolbeault, David Krieg, Mario Ullrich}
\date{Source: arXiv:2602.18719v1 (CC BY 4.0)}
\begin{document}
\maketitle

\noindent\emph{Source and licence.} Constructive discretization and approximation in reproducing kernel Hilbert spaces. Version arXiv:2602.18719v1 (CC BY 4.0), distributed under the Creative Commons Attribution 4.0 International licence, as recorded in the arXiv licence metadata for this exact version and shown on the abstract page https://arxiv.org/abs/2602.18719v1. Full rights notice: licenses/arxiv-2602.18719-CC-BY-4.0.txt.\par\smallskip

\noindent\textit{Editorial scope.} Counted unit: the Coro whose source label is \texttt{coro:main} in the article (source0.tex lines 284--306, source label \texttt{coro:main}), delivered together with the article's own complete original proof of it (source0.tex lines 308--334, one \texttt{proof} environment of 27 lines). No mathematics is written, repaired or re-derived: statement and proof are the article's own text line for line. Required context retained uncounted: thm:main (lines 144--168); cond:continuous (lines 268--270); cond:injective (lines 272--275); eq:finite-trace (lines 277--279). Every definition, display and label of those lines is retained unchanged. Omitted from this package and not redistributed: 1--143, 169--267, 271--271, 276--276, 280--283, 307--307, 335--1919. Nothing inside a retained excerpt is omitted or altered: every retained body is delivered exactly as the article's lines are written, so no omission receipt of any kind accompanies this package. Citations: the retained excerpts carry no citation command at all, so no bibliography entry of the article is transcribed and no reference is redistributed. Reference adaptation: none was needed and none was made; every reference inside the delivered unit resolves inside it, so no unresolved reference or citation is produced. Printed slips kept literal: no printed word, symbol, hypothesis or number of the counted statement or of its proof was altered. Layout and class: the article's own class is \texttt{article} with packages including inputenc, amsmath, amsthm, amsfonts, amssymb, stmaryrd, graphicx, xcolor, paralist, soul, framed, setspace, tikz-cd, bbm; the wrapper is the lane's pinned \texttt{amsart} editorial wrapper at 10pt on A4, which restates the theorem-like environments the retained text needs and adds the article's own macro definitions that the retained text needs (\texttt{\textbackslash C}, \texttt{\textbackslash II}, \texttt{\textbackslash Tr}), with extra wrapper packages (bm, bbm, dsfont, xspace, cancel, mathtools). The delivered numbering is the unit's own. Assets: no retained span contains a figure environment, a graphics-inclusion command, a PDF-page inclusion, a table or a TikZ picture; no image file is copied into this package, the package declares no asset dependency and no image file is redistributed with it. Licence: version arXiv:2602.18719v1 of this article is distributed under the Creative Commons Attribution 4.0 International licence, as recorded in the arXiv licence metadata for this exact version and shown on the abstract page https://arxiv.org/abs/2602.18719v1. A full rights notice is delivered at licenses/arxiv-2602.18719-CC-BY-4.0.txt. Characters: every retained excerpt is pure ASCII, so the delivered engine, XeLaTeX with its default Latin Modern text font, reports no missing character.
\begin{theorem}
\label{thm:main}
Let $(D,\mu)$ be a measure space, and let $a=(a_1,\hdots,a_m)^\top$
and $b= (b_k)_{k\in \mathbb I}^\top$ be families of square-integrable functions on $D$,
where $\mathbb I$ is at most countable.
Assume that $I:=\int_D a(x)a(x)^*d\mu(x)$ and $J:=\int_D b(x)b(x)^*d\mu(x)$ are positive definite, and that $J$
has finite trace and effective dimension
\[
\Tr(J) \,=\, \sum_{k\in \mathbb I} \|b_k\|_{L_2(D,\mu)}^2 \,<\,\infty
\quad\text{and}\quad M:=\frac{\Tr(J)}{\lambda_{\max}(J)}.
\]
Then, for any
$n\geq m$, there exist points $x_1,\dots,x_n\in D$ and weights ${w_1,\dots,w_n>0}$ such that
\[
\lambda_{\min}\left( \sum_{i=1}^{n} w_i\, 
a(x_i)a(x_i)^* \right) \geq
\left(1-\sqrt{\frac{m-1}{n}}\right)^2
\lambda_{\min}(I)
\]
and
\[
\lambda_{\max}\left(\sum_{i=1}^{n} w_i\, b(x_i)b(x_i)^* \right) \leq \left(1+\sqrt{\frac{M-1}{n}}\right)^2
\lambda_{\max}(J).
\]
\end{theorem}

\begin{equation}\label{cond:continuous}
\lambda:=\sup_{f\in H} \frac{\Vert f \Vert_2^2}{\Vert f \Vert_H^2}<\infty
\end{equation}

\begin{equation}\label{cond:injective}
 \Vert f \Vert_2 \ne 0
 \quad \text{for all} \quad f\in H \setminus \{0\},
\end{equation}

\begin{equation}\label{eq:finite-trace}
 \Tr(K) := \int K(x,x)\,d\mu(x) < \infty.
\end{equation}

\begin{coro} \label{coro:main}
Let $(D,\mu)$ be a measure space, 
$V_m\subset L_2(D,\mu)$ be an $m$-dimensional space of functions and let $H$ be 
a reproducing kernel Hilbert space satisfying conditions \eqref{cond:continuous}, \eqref{cond:injective} and \eqref{eq:finite-trace}. 
For $n\geq m$, there exist points $x_1,\dots,x_n\in D$ and weights ${w_1,\dots,w_n>0}$ such that
\begin{equation}
\label{eq-low-coro-main}
\left(1-\sqrt{\frac{m-1}{n}}\right)\,\|f\|_2
\,\leq\,
\sqrt{\sum_{i=1}^n w_i|f(x_i)|^2}
\qquad \text{for all } f \in V_m
\end{equation}
and 
\begin{equation}
\label{eq-up-coro-main}
\sqrt{\sum_{i=1}^n w_i|g(x_i)|^2}
\,\leq\,
\left(1 +\sqrt{\frac{M-1}{n}} \right)\, \sqrt\lambda\,
\left\| g \right\|_H
\qquad \text{for all } g \in H,
\end{equation}
where $M=\Tr(K)/\lambda$ is the ``effective dimension'' of $H$ in $L_2(D,\mu)$.
\end{coro}

\begin{proof}
Let $a=(a_1,\hdots,a_m)^\top$ be an $L_2$-orthonormal basis of~$V_m$,
and $b=(b_1,b_2,\hdots)^\top$ be the basis of left singular vectors of the embedding $ H\hookrightarrow L_2$.
Then $I=\int a\,a^*d\mu$ is the identity, and
$J=\int b\,b^*d\mu$
is diagonal with trace $\Tr(J) =\Tr(K)$ and largest entry $\lambda_{\max}(J) = \lambda$,
since $b$ is orthonormal in $H$ and orthogonal in $L_2$.

Let $x_i$ and $w_i$ be the points and weights from Theorem~\ref{thm:main}, and denote
\[
G=\sum_{i=1}^{n} w_ia(x_i)a(x_i)^*\quad\text{and}\quad
\Gamma = \sum_{i=1}^{n} w_ib(x_i)b(x_i)^*.
\]
Then, any function $f\in V_m$ can be written as $f=\sum_{k=1}^{m} c_k a_k$ 
with $c\in \C^m$. Hence,
\[
 \sum_{i=1}^n w_i|f(x_i)|^2 = c\,G\,c^*
\geq \lambda_{\min}(G)\,|c|^2
 =\lambda_{\min}(G)\,\Vert f\Vert_2^2.
\]
In the same way, decomposing any $g\in H$ as $g=\sum_{k\in \mathbb I} c_k b_k$ with $c\in \ell_2(\II)$, we get\vspace{-4mm}
\[
 \sum_{i=1}^n w_i|g(x_i)|^2 = c\,\Gamma\,c^*
 \leq \lambda_{\max}(\Gamma)\,\Vert c \Vert_{\ell_2(\II)}^2
 =\lambda_{\max}(\Gamma)\,\Vert g \Vert_{H}^2.\qedhere
\]
\end{proof}

\end{document}
